\documentclass[10pt]{article}
\usepackage[margin=1in]{geometry}
\usepackage{amsmath}

\begin{document}

\noindent\fbox{%
\begin{minipage}{0.96\textwidth}
\textbf{Analysis of acid--base metadynamics trajectories.}
For each acid--base metadynamics trajectory, a $3\times3$ diagnostic summary was constructed by aligning the trajectory, Mulliken-population, bias-potential, and free-energy records according to their reported simulation times. All intermolecular distances and hydrogen-bond geometries were evaluated using the periodic minimum-image convention. The positive protonic defect was assigned in every frame to the solvent oxygen with the largest Mulliken charge within the prescribed defect interval, $-0.625\leq q_{\mathrm O}\leq-0.525\ e$. In panel (a), the one-dimensional free-energy profile $F(s)$ was taken from the final internally reconstructed free-energy block available at or before $t_{\mathrm{diff}}+1.75$~ps, converted from hartree to kcal~mol$^{-1}$, and shifted such that its minimum was zero. Here, $s$ is the on-the-fly coordination variable between the designated solute nitrogen and the biased hydrogen. The simulations employed conventional (non-well-tempered) metadynamics, in which Gaussian bias potentials deposited along $s$ are accumulated to give $V_{\mathrm{bias}}(s,t)$. As the bias progressively fills the underlying free-energy wells, the traditional metadynamics estimator used internally by DCDFTBMD reconstructs the free-energy surface as $F(s,t)\simeq-V_{\mathrm{bias}}(s,t)+C(t)$, where $C(t)$ is an arbitrary time-dependent additive constant; thus, the reported profile is the finite-time accumulated-bias estimate and not a separate post-processing histogram or Tiwary--Parrinello reweighting. Panel (b) reports the instantaneous hydrogen-bonded wire connecting the designated nitrogen to the defect-bearing oxygen. A hydrogen bond D--H$\cdots$A was assigned when $d(\mathrm{D-H})\leq1.3$~\AA, $d(\mathrm{H}\cdots\mathrm{A})\leq2.5$~\AA, and $\angle\mathrm{D-H\cdots A}\geq135^{\circ}$; breadth-first search then identified the shortest connected path containing no more than four intervening water molecules. The plotted wire state is the number of bridging waters, with 0 denoting direct N--defect contact and $-1$ denoting a disconnected path, and the connected and disconnected probabilities were evaluated over $t_{\mathrm{diff}}<t\leq t_{\mathrm{diff}}+1.75$~ps. Panel (c) shows both the N--defect distance and the distance from the defect to the closest solute heavy atom; N--defect separations were partitioned into the predefined solvation layers 2.0--3.5, 3.5--5.5, 5.5--7.5, and 7.5--9.5~\AA. Panel (d) presents the biased coordination time series $s(t)$, separating the reaction and diffusive intervals at $t_{\mathrm{diff}}$. Unless imposed manually, $t_{\mathrm{diff}}$ was assigned after detection of a distance-separated state that remained at $s\leq0.05$ for the specified persistence interval and was followed by sustained recovery to $s\geq0.20$; the separation threshold and persistence interval were treated as configurable analysis parameters. Panel (e) characterizes water-wire conductivity using $\xi=N_b^{-1}\sum_{i=1}^{N_b}\sigma_i$, where the $N_b$ hydrogen bonds are ordered from the defect toward the nitrogen and $\sigma_i=+1$ ($-1$) when the donor lies on the defect-facing (nitrogen-facing) side of bond $i$. The displayed $P(\xi)$ is the normalized discrete distribution in the post-diffusion interval; neighboring bond orientations are labeled P, D, L, and R for $(+,+)$, $(+,-)$, $(-,+)$, and $(-,-)$, respectively. Panel (f) displays the Mulliken charges of all solvent oxygen atoms, calculated by summing their reported $s$- and $p$-orbital contributions; vertical markers indicate changes in the identity of the assigned defect oxygen and therefore candidate Grotthuss proton-transfer events. Panel (g) gives the time-dependent acidity as $\mathrm{p}K_a(t)=\Delta F(t)/(RT\ln10)$, with $T=313.15$~K by default and $\Delta F=F(s\approx0)-F(s\approx1)$. At every saved FES time, the two basin minima were located within $\pm0.1$ of their nominal coordination values. The reported $\mathrm{p}K_a$ and uncertainty are the mean and population standard deviation of ten linearly interpolated, equally spaced values between $t_{\mathrm{diff}}$ and $t_{\mathrm{diff}}+1.75$~ps. Panel (h) reports the loopiness of each connected wire as $\lambda=1-R/L$, where $L$ is the sum of the minimum-image lengths of all consecutive path segments and $R$ is the magnitude of their vector sum; hence, $\lambda=0$ describes a straight path and increasing $\lambda$ indicates greater tortuosity. The inset is a boundary-reflection-corrected Gaussian kernel-density estimate of $P(\lambda)$ over the post-diffusion interval. Finally, panel (i) reports the hydrogen-bond coordination $m$ of the defect-bearing oxygen, defined as its degree in the same hydrogen-bond graph---the number of unique solvent-oxygen or designated-nitrogen neighbors satisfying the above geometric criteria---with the inset giving the corresponding normalized discrete distribution $P(m)$ over $t_{\mathrm{diff}}<t\leq t_{\mathrm{diff}}+1.75$~ps.
\end{minipage}%
}

\end{document}
